
\begin{abstract}
As language models gain autonomy, they increasingly act as \emph{self-interested}
agents on behalf of distinct principals rather than as cells of one cooperative
organism. Such agents need institutions: shared, enforceable, disputable rules.
We argue that the substrate for these institutions should be \emph{alignment-by-institution}---an
explicit, inspectable, defeasible normative layer \emph{over} capable reasoners---rather
than only alignment-by-internals (steering vectors, interpretability). We
instantiate this with a defeasible deontic logic engine (Governatori-style proof
theory) implemented in Lean~4 that compiles normative text into theories and
computes obligations, permissions, prohibitions, compensatory remedies, violations,
and \emph{unresolved} conflicts. We report a pilot evaluation and a sequence of
small, adversarially-designed experiments. Three findings stand out. (i) On a
penal-law verdict benchmark the engine and an LLM alone are near-tied on the
\emph{guilty/not} question (100\% vs.\ 96\%), but diverge sharply on \emph{which}
norm applies (offence-ID 83\% vs.\ 33\%; retrieval-augmented 0\%)---the metric an
institution actually needs. (ii) The engine's labels are corroborated by
\emph{independent} annotators from two different model families (gpt-4.1 89\%,
DeepSeek~V3 83\%), addressing the charge of circular self-grading; and in an
illustrative contract-breach scenario two self-interested LLM
agents reach \emph{opposite} verdicts while the engine yields one verdict both can
pre-commit to. (iii) The real bottleneck is the text$\rightarrow$logic
front-end: unscaffolded single-shot formalization is unreliable (0/3), but the
failure \emph{decomposes} into a cheap syntax layer (solved by a grammar spec)
and a semantics layer that, given operator \emph{definitions}, generalizes
correctly even to clause shapes never demonstrated (out-of-distribution
contrary-to-duty: pass); a first multi-clause test, however, surfaces a
compound-effect failure mode, leaving document-scale fidelity open. Throughout, a defining property holds: malformed
formalizations \emph{fail loudly} (rejected by the loader) or are legible on
inspection---a verification asset weight-space methods lack. (iv) We then stress this
property adversarially: a self-interested scribe can poison an encoding so it loads yet
escapes an obligation. An independent cross-family reviewer catches a blatant poison but
only \textbf{1 of 3} subtle structural ones (a deleted compensation chain, an inverted
superiority, a flipped condition)---so inspectability is necessary but \emph{not} sufficient.
Grounding the reviewer in the engine's \emph{computed} verdict on an agreed probe (engine-assisted
audit) restores detection to \textbf{3 of 3}, at some cost to precision---locating the defence in
differential testing, not free-form review. We are explicit that
several studies are single illustrative scenarios, and that the natural-language$\rightarrow$logic
formalization itself is studied more thoroughly by concurrent work; our novel emphasis is the
institutional and adversarial framing.
\end{abstract}

\section{Introduction}
Most multi-agent LLM systems today behave as a single organism optimizing one
shared objective. But the more valuable---and harder---regime is a \emph{society}
of agents with different, sometimes conflicting, goals that nonetheless benefit
from exchange, specialization, and contracting. Human societies enable this with
institutions and the rule of law: machinery to resolve disputes, apply rules, and
identify which norms govern a case. We ask what the equivalent machinery is for AI
agents, and argue for an \emph{engineered} normative layer rather than an organically
evolved one.

The dominant alignment toolkit operates on a single model's internals: steering
vectors, mechanistic interpretability, constitutional training. These are valuable
but address the wrong locus for \emph{inter-agent} norms. ``Agent~$A$ breached its
delivery obligation to Agent~$B$'' is not a fact in any one network's activations;
it is a relational fact about a joint history plus a rule. We therefore explore
\emph{alignment-by-institution}: an explicit, inspectable, machine-checkable
normative layer that sits \emph{over} capable reasoners and that self-interested
parties can agree to in advance.

\paragraph{Contributions.}
Using deontic norms to govern agents is an old idea (\S\ref{sec:related}); we do not claim the
formalization task itself as new (cf.\ \citet{horner2025robust}). Our contributions are:
(1) an \emph{institutional} reframing for \emph{self-interested} agents---the engine as a
\emph{commitment device} two parties pre-commit to, evidenced by a contract-breach scenario where
self-interested LLM agents diverge but the engine does not (E3); (2) an \emph{adversarial}
threat model and a first result (E12): a self-interested scribe can poison the \texttt{.ddl} so it
loads yet escapes a penalty, but an independent cross-family reviewer catches it---auditability as
a \emph{security} property, which prior NorMAS and LLM-formalization work does not study;
(3) a backward-generated evaluation isolating \emph{norm identification} from \emph{verdict},
showing the engine's measurable edge is the former; and (4) a small replication that decomposes the
text$\rightarrow$logic bottleneck into a solved syntax layer and a narrower semantics layer.
The engine itself (Lean~4: Hohfeldian bearers, compensatory chains, \texttt{[JUDGE]} conflicts)
is an implementation of an existing calculus, not a contribution.

\paragraph{Methodological note.}
This paper was developed through a structured adversarial debate between a
proponent persona and a skeptic persona favoring non-symbolic methods. Each
contested claim was settled, where possible, by a small experiment rather than
argument. Conceded points are reported as such; the skeptic won at least one round
(the formalization bottleneck), and the paper reflects that.

\section{The Engine}
We use defeasible deontic logic in the proof-theoretic formulation of Governatori
(Figure~\ref{fig:pipeline} shows the full pipeline).
A theory is a set of labelled rules with strengths (strict \texttt{->}, defeasible
\texttt{=>}, defeater \texttt{\textasciitilde>}), a superiority relation between
rule labels, and a separation between constitutive and prescriptive ($\dots O$)
rules. Conclusions may be \emph{compensatory chains} $a * b * c$: $O(a)$, and on
violation of $a$, $O(b)$, etc.\ (contrary-to-duty). The engine computes a per-bearer
fixed point and reports, per literal, whether it is definitely provable in each
deontic mode. Two outputs matter institutionally: \texttt{[JUDGE]} flags conflicts
the rules under-determine (the system refuses to guess the law), and a
\emph{non-compensable violation} flags a breach with no remaining remedy.

\paragraph{Grounding.} Atom names are opaque, so every used atom must carry a
natural-language description; the loader \emph{refuses} undescribed atoms. This
makes a theory's meaning legible and a wrong rule diagnosable by reading it---a
property we exploit below.

\begin{figure}[h]\centering
\begin{tikzpicture}[
  font=\small,
  box/.style={draw, rounded corners, align=center, inner sep=4pt, minimum height=9mm},
  out/.style={draw, align=center, inner sep=3pt, fill=black!4},
  >={Stealth[]}]
\node[box] (text) {Normative\\text};
\node[box, right=10mm of text] (llm) {LLM\\formalizer\\{\scriptsize(grammar spec)}};
\node[box, right=10mm of llm] (ddl) {\texttt{.ddl} theory\\{\scriptsize grounded atoms +\\defeasible rules}};
\node[box, right=10mm of ddl] (eng) {Deontic engine\\{\scriptsize per-bearer\\fixed point}};
\node[out, right=10mm of eng] (res) {$O/P/F$, comp.\\chains; \texttt{[JUDGE]};\\non-comp.\ violation};
\draw[->] (text) -- (llm);
\draw[->] (llm) -- (ddl);
\draw[->] (ddl) -- (eng);
\draw[->] (eng) -- (res);
\node[align=center, below=6mm of ddl, text width=42mm] (load)
  {\scriptsize\emph{loader rejects undescribed atoms / malformed rules} $\Rightarrow$ \textbf{loud failure}};
\draw[->, dashed] (ddl) -- (load);
\end{tikzpicture}
\caption{Pipeline. The LLM is the (fallible) \emph{scribe}; the engine is the
deterministic reasoner. The auditability property (dashed) is that bad theories are
rejected at load or are legible in the \texttt{.ddl}, never silently wrong.}
\label{fig:pipeline}
\end{figure}

\section{Experiments}
Experiments use \texttt{gpt-4.1} where an LLM is required, unless noted (E9 uses
DeepSeek~V3). Datasets and harness are small (pilot scale); we are explicit about
$n$ and treat results as directional.

\subsection{Norm identification vs.\ verdict (E1)}
We build a penal-law benchmark by \emph{backward generation}: the engine over an
encoded slice of the Italian penal code is the gold oracle; we pick an atom
configuration, read off the engine's verdict as the label, then dress it as an
Italian narrative with no element-phrasing or leakage, including minimal pairs that
flip one decisive element. We compare three arms: LLM-only, retrieval-augmented
(RAG), and LLM+engine (deontic).

\begin{table}[h]\centering
\begin{tabular}{lccc}
\toprule
Arm & Verdict acc. & Offence-ID & Pair-acc.\\
\midrule
deontic   & \textbf{100\%} [88--100] & \textbf{83\%} (15/18) [61--94] & 100\% (11/11) [74--100]\\
llm\_only  & 96\% [82--99]  & 33\% (6/18) [16--56]  & 91\% (10/11) [62--98]\\
rag       & 96\% [82--99]  & 0\% (0/18) [0--18]    & 91\% (10/11) [62--98]\\
\bottomrule
\end{tabular}
\caption{Pilot eval, $n{=}28$ (offence-ID over the 18 offence items; 11 minimal pairs).
Brackets are Wilson 95\% confidence intervals. Grounding accuracy (deontic) is 94\%
(92/98 atoms). RAG \emph{underperforms} the bare LLM on offence-ID, evidence against
``just retrieve the statute.''}
\label{tab:e1}
\end{table}

Three reads, with the CIs disciplining each. (i) The \emph{verdict} gap is not
distinguishable---all three intervals overlap heavily---consistent with the verdict
question (``was a crime committed'') being nearly guessable from narrative tone. (ii)
The \emph{offence-ID} gap \emph{is} real: deontic $83\%$ [61--94] vs.\ llm\_only $33\%$
[16--56] have non-overlapping intervals even at $n{=}18$, and RAG at $0\%$ [0--18] is
cleanly separated. This is the institutional question---\emph{which} offence (lex
specialis) governs, hence the penalty and remedy. A caveat we make explicit (and
revisit in \S\ref{sec:discussion}): a manual audit of the llm\_only misses finds
$\sim$80\% are the \emph{correct} offence named non-canonically (e.g.\ the element
atom, or an invented synonym), so this gap is better read as the engine performing
\emph{canonicalization} than as superior legal reasoning. (iii) We \emph{decline} to claim the
minimal-pair difference ($100\%$ vs.\ $91\%$): at $n{=}11$ the intervals overlap and the
difference is not statistically supported, a point we flag rather than oversell.

\subsection{Breaking circularity (E2)}
Objection: the engine grades the LLM against gold the engine itself produced. We
test this with an \emph{independent} annotator: \texttt{gpt-4.1} given the full
cross-offence catalogue and room to deliberate, never touching the engine. It agrees
with the engine's offence-ID gold on \textbf{16/18 (89\%)}. The two disagreements are
illuminating, not damning: one is a genuine interpretive fork (an aggravation the
config did not set)---exactly the \texttt{[JUDGE]} case; the other is a premeditation
$\rightarrow$ life-sentence escalation the \emph{independent model also missed} and
only the engine's enforced lex-specialis caught.

A sharper version of the objection is that this annotator is \emph{itself} \texttt{gpt-4.1},
the same model family as the arm under test, so agreement could reflect shared priors rather
than legal correctness. We therefore repeat the check with a \emph{different} model family,
DeepSeek~V3, which agrees with the engine gold on \textbf{15/18 (83\%)}. Two independent model
families thus corroborate the gold at 83--89\%, which addresses the same-model concern. Notably,
the premeditation$\rightarrow$\texttt{Ergastolo} escalation is missed by \emph{both} families and
caught only by the engine---a model-independent instance of enforced lex specialis; DeepSeek's
other two misses are genuine legal errors, consistent with a weaker legal reasoner that is still
83\% aligned. Finally, single-pass offence-ID was 33\% but deliberation-with-canonical-vocabulary
reaches 83--89\%: the gap is closed by supplying canonical vocabulary and forcing
element-by-element reasoning---which is what the engine institutionalizes.

\subsection{The commitment device (E3)}
We encode a two-party contract: the seller owes on-time delivery with a contrary-to-duty
late penalty (\texttt{=>O@Seller DeliverOnTime * PayPenalty}); the buyer owes payment
on delivery. In the ``delivered late'' scenario the engine deterministically derives
$O(\text{PayPenalty})$. We then ask \texttt{gpt-4.1} to argue as each \emph{self-interested}
party. The seller's counsel argues the penalty is \emph{not} owed (invoking
\emph{exceptio non adimpleti contractus}---a real doctrine, but a cross-default the
contract does not encode); the buyer's counsel argues it is; a neutral arbiter and the
engine agree it is. The parties \emph{diverge}; the engine is the contract-as-written
referee both can pre-commit to before knowing their side. The seller's defence also
localizes the next problem: whether the \emph{exceptio} applies is a \emph{formalization}
question the theory settles up front.

\subsection{The formalization bottleneck, decomposed (E4--E8)}
We probe text$\rightarrow$logic fidelity end-to-end: prose in, \texttt{.ddl} out, loaded
in the engine and scored against a gold verdict.

\begin{itemize}
\item \textbf{E4 (scaffolded, simple):} given atom-name hints and the CTD pattern,
the model round-trips a clause to valid \texttt{.ddl} reproducing the gold verdict. Feasible---but
the scaffolding may be doing the work.
\item \textbf{E5 (unscaffolded, textbook clause):} a prohibition with a gated permission.
The model emits fluent-looking but \emph{semantically wrong} logic: it reifies permission
into an obligated proposition $O(\text{permitted})$ instead of using the permission operator
on the act, and inverts the superiority direction. Loads, but wrong.
\item \textbf{E6 (replication, 2 fresh clauses):} the E5 error does not recur (both use
\texttt{\textasciitilde>} correctly) but both \emph{fail to load} on new errors---``\&'' for
conjunction where the grammar wants a comma, and superiority over labels never assigned.
Unscaffolded tally: \textbf{0/3 correct}.
\item \textbf{E7 (layer isolation):} re-running the E6 clauses with a syntax spec plus one
worked permission-exception example yields \textbf{2/2}, syntactically and semantically
correct. The syntax layer is scaffolding-solvable; the semantics error vanishes given a
template.
\item \textbf{E8 (out-of-distribution):} the grammar spec \emph{defines} the CTD ``$*$''
operator but the worked example never shows it; the test clause (late rent $\rightarrow$
rent + penalty) requires it. \textbf{Pass:} the model deploys ``$*$'' correctly and the
engine derives the gold verdict. Generalization from a \emph{definition}, not a demonstration.
Caveat: the encoding is non-minimal (a redundant parallel rule and a spurious superiority)---correct
but inelegant, and legibly so.
\item \textbf{E9 (second model, DeepSeek~V3 via OpenRouter):} re-running E8 on a non-OpenAI
model, the \emph{semantic} generalization replicates---it deploys ``$*$'' correctly from the
definition---but the file \emph{fails to load} on a surface lexical-consistency bug (a misspelled
atom name in a rule vs.\ its declaration), which the engine flags exactly. Both models independently
add the \emph{same} redundant rule and spurious superiority. So semantics-generalization is
cross-model; \emph{surface-fidelity reliability is model-dependent} (the stronger model loads, the
weaker trips on lexical consistency); and the minimality issue is model-independent.
\end{itemize}

\textbf{E10 (multi-clause scale).} A five-clause contract with \emph{interacting}
defeasibility (non-payment derogates the service duty; a security incident triggers a
notify$*$compensate chain; force majeure derogates). The model produces 4/5 clauses
correctly---including the contrary-to-duty chain and \emph{both} superiority relations---and,
notably, keeps the atom vocabulary \emph{consistent} across eight atoms and five rules, so
cross-clause coreference does not break. The failure is localized to the single clause
requiring a \emph{compound} effect (permit suspension \emph{and} release the service
obligation): the model crams two conclusions into one rule, and the theory fails to load.
A new failure mode---multi-effect clauses (one clause maps to $\geq 2$ rules)---appears under
composition, but is caught by the loader, and the global structure (superiorities, bearers,
shared vocabulary) is right. (We also observed the model intermittently emit a \emph{filename}
in place of theory content on this larger task: schema adherence itself degrades with prompt
complexity.)

\paragraph{Arc.} Unscaffolded 0/3 $\rightarrow$ syntax-spec + isomorphic example 2/2
$\rightarrow$ spec-defined-operator OOD 1/1 $\rightarrow$ multi-clause 4/5 clauses (one
compound-effect clause fails). The bottleneck is real but \emph{decomposes}:
a cheap syntax layer (grammar spec / constrained decoding) and a semantics layer that,
given operator definitions, is more tractable than the raw 0/3 suggested. Across all
failures the system \emph{fails loud or legible}: malformed theories are rejected by the
loader; wrong-but-loadable ones are diagnosable by reading the \texttt{.ddl}. This
auditability is the structural advantage over weight-space encodings of the same norm.

\subsection{Adversarial formalization (E12)}
Every experiment above assumes a cooperative scribe. In a society of self-interested agents
the party who writes the \texttt{.ddl} has motive to cheat. We test this: the seller's agent is
asked to encode the late-delivery contract so it \emph{loads} yet the penalty never applies.
The attack \emph{succeeds}---the model smuggles a fabricated precondition (penalty owed only if
delivery is ``certified by an independent moderator,'' a condition absent from the contract) into
the contrary-to-duty rule, so the late-delivery scenario yields no $O(\text{Penalty})$. But the
\emph{defence also succeeds}: an independent model from a different family (DeepSeek~V3), given the
prose contract and the poisoned \texttt{.ddl}, flags the smuggled condition precisely. The
auditability property thus extends to the adversarial setting for the reason it holds elsewhere:
the manipulation is a \emph{visible extra atom} in an inspectable artifact, not a hidden weight
perturbation. The institutional reading is that one does not trust the scribe; the counterparty's
(or a neutral) agent audits the \texttt{.ddl} against the agreed prose.

This optimism does not survive harder attacks. In E13 we hand-author three \emph{subtle}
poisons that add \emph{no new vocabulary}---a silently deleted compensation chain (the
\texttt{* PayPenalty} simply dropped), an inverted superiority (a no-penalty rule made to defeat
the contrary-to-duty rule), and a condition flip (penalty owed only if delivery is \emph{on
time})---each of which loads and lets the late-delivery penalty escape. The same independent
reviewer now catches only \textbf{1 of 3}: it flags the inverted superiority (which adds a
visible rule) but misses both the deletion and the condition flip; on the latter it even restates
the contract correctly and then declares the faithless encoding faithful. The lesson is sharp:
\emph{inspectability is necessary but not sufficient}. A representation being legible does not mean
a naive single-pass reviewer detects a structural manipulation in it. Robust adversarial audit is
therefore the central open problem for this approach---and the natural direction is
\emph{engine-assisted} audit (e.g.\ differential testing: have the engine compute the theory's
verdict on agreed probe scenarios and confront the reviewer with ``this encoding yields \emph{no}
penalty on late delivery''), turning shallow prose-vs-code review into grounded review.

We test that fix directly (E14): we run the engine on the agreed late-delivery probe and hand the
reviewer its \emph{computed verdict} as a neutral fact (``$O(\text{PayPenalty})$ holds: yes/no'')
alongside the prose and the \texttt{.ddl}. Detection of the three structural attacks rises from
$1/3$ to \textbf{$3/3$}, stable across three runs---the reviewer now connects ``the contract
requires the penalty when late'' with ``this theory yields no penalty when late.'' The cost is a
precision one: in $1$ of $3$ runs the reviewer also flagged the \emph{faithful} control (a false
positive), so grounding raises recall reliably but needs calibration. The engine's role in
adversarial audit is thus not to out-reason the LLM but to \emph{ground} it: a deterministic
behavioural verdict on an agreed probe converts a manipulation invisible in the syntax into a
discrepancy the reviewer can see. We report E12--E14 as single illustrative attacker/reviewer
studies (one model each, small per-cell $n$), not distributions; the strength is the
\emph{direction}---inspectability plus engine-grounded differential audit---not the point estimates.

\section{Discussion: what is settled, what is open}\label{sec:discussion}
\textbf{Settled (on this evidence).} The engine's defensible value is not ``better legal
reasoning'' than the LLM---on near-saturated verdicts they tie, and $\sim$80\% of the LLM's
offence-ID ``errors'' in early analysis were correct reasoning under non-canonical atom names.
Its value is \emph{canonicalization} (binding correct-but-idiolectal reasoning to one shared,
machine-checkable vocabulary), \emph{enforced lex specialis} (escalations LLMs silently drop),
\emph{conflict surfacing} (\texttt{[JUDGE]}), \emph{commitment-device adjudication} for
self-interested parties, and \emph{auditability}. These make it a plausible \emph{protocol layer}
for agent institutions, orthogonal to---not competing with---interpretability and steering.

\textbf{Open (the central problem).} Formalization fidelity at scale. Our evidence is
clause-level and small-$n$. A first multi-clause probe (E10) is encouraging on global structure
(superiorities, bearers, consistent shared vocabulary across five rules) but exposes a
compound-effect failure mode (one clause $\mapsto \geq 2$ rules) and degraded schema adherence.
Still untested: contracts with deep cross-references and nested defeasibility at document scale;
the minimality of encodings under composition; broad model coverage; adversarial clauses designed
to induce \emph{silent} omissions (we have seen only loud ones so far). The decisive future
benchmark contains structurally novel, multi-effect, and document-scale clauses, scored
end-to-end against engine-gold verdicts.

\section{Limitations}
$n$ is pilot-scale (28 eval items; 1--3 clauses per formalization condition). The eval verdict/offence arms (E1)
use \texttt{gpt-4.1}; the independent-annotator check (E2) and the formalization OOD test (E9)
are run on both \texttt{gpt-4.1} and DeepSeek~V3, but broader model coverage is still thin.
Quantitative claims remain directional pending larger samples. The penal-law gold is the
engine's own output (mitigated, not eliminated, by the independent-annotator check). The
commitment-device and formalization studies are illustrative single scenarios, not distributions.
We deliberately report these rather than over-claim.

\section{Related Work}\label{sec:related}
\paragraph{Defeasible deontic logic.} We build directly on the proof-theoretic tradition of
defeasible deontic logic, in particular reasoning about violations and contrary-to-duty
obligations via a Gentzen-style system~\citep{governatori2006violations} and recent treatments
of compensatory chains that avoid pragmatic oddity~\citep{governatori2024oddity}. Our engine
is an implementation of this calculus; our contribution is not the logic but its use as
institutional infrastructure and an empirical study of the natural-language front-end. We model
directed (party-bearing) obligations using Hohfeld's analysis of jural
relations~\citep{hohfeld1913}.

\paragraph{Alignment by training vs.\ by institution.} Constitutional AI~\citep{bai2022constitutional}
trains a fixed set of principles \emph{into} a model---a \emph{civil-law} stance, where the rules
are compiled once and the model is trusted to apply them. We explore the complementary
\emph{common-law} stance: an explicit, external, case-extensible normative layer that adjudicates
\emph{between} agents and surfaces, rather than silently resolves, genuine conflicts. This is
orthogonal to weight-space alignment: a normative fact about two agents is not located in one
model's parameters.

\paragraph{Scalable oversight.} AI safety via debate~\citep{irving2018debate} and
weak-to-strong generalization~\citep{burns2023weak} scale supervision with model capability.
We treat these as complementary, not competing: debate produces a \emph{winner}; an institution
produces a \emph{contract both parties agreed to in advance}. For self-interested agents the latter,
not the former, is the trust primitive.

\paragraph{Normative MAS and electronic institutions.} Using deontic norms to govern agent
societies is not new: normative multi-agent systems~\citep{boella2006normative} and electronic
institutions with enforcement middleware~\citep{esteva2004ameli} pursued exactly this---obligations,
permissions, and prohibitions coordinating self-interested agents in electronic commerce---two
decades ago. The classical setting, however, assumed agents are \emph{purpose-built software} and
norms are \emph{hand-authored} by an institution designer; the open problem of feeding messy
natural language into the formal layer was left aside because there was no general text understander.
Our premise is that this gap is now the crux: the agents are general LLM reasoners, and the LLM is
also the formalization front-end. We therefore inherit the NorMAS goals but relocate the hard
problem to the text$\rightarrow$logic bridge and its adversarial robustness.

\paragraph{LLM-to-deontic-logic formalization.} Closest to our formalization experiments is
\citet{horner2025robust}, who formalize legal text into defeasible deontic logic with a multi-stage
LLM pipeline and a dedicated refinement step, evaluated on a real statute across model
configurations; related auto-formalization-with-verification efforts target regulation such as
GDPR~\citep{nguyen2026gdpr}. \textbf{We do not claim the formalization task as novel: \citet{horner2025robust}
study it more thoroughly than our small probes (E4--E11), and our results are best read as a
small replication and decomposition (syntax vs.\ semantics layers).} Our distinct contributions are
elsewhere: (i) the \emph{institutional} framing for self-interested agents---the engine as a
\emph{commitment device} both parties pre-commit to (E3); (ii) the \emph{adversarial} threat model
(E12)---a self-interested scribe poisons the encoding, and an independent cross-family reviewer
catches it---which neither the NorMAS nor the LLM-formalization literature addresses; and (iii) the
separation of \emph{norm identification} from \emph{verdict} and the auditability-as-security
argument. This connects to autoformalization in mathematics~\citep{wu2022autoformalization}, with
enforceable norms rather than theorems as the target, and to LegalBench~\citep{guha2023legalbench},
which measures legal \emph{reasoning} where we measure \emph{formalization}.

\section{Conclusion}
An explicit, auditable, defeasible normative layer is a credible substrate for institutions
among self-interested AI agents. The calculus is sound and the adjudication properties are
real; the binding constraint is the natural-language$\rightarrow$logic front-end, which our
experiments show is neither trivial nor hopeless---it decomposes into a solved syntax layer and
a tractable, definition-driven semantics layer, with scale and minimality as the genuine
remaining risks. Building and stress-testing a structurally-diverse formalization benchmark is
the next step.
